% ===================================================================
%  جدول نمبر 8 — فصل کی کٹائی۔
%
%  Traduction, faite en 2026, en ourdou d'aujourd'hui, sur l'ido de
%  texto/io. Regles de la colonne : voir 10-jadval-01.tex. Deux
%  scenes, deux numerotations : les cles portent c1 et c2. Ecrit avec
%  outils/ordo.py sous les yeux. Le fac-simile ido compose « 13) »
%  sans sa parenthese ouvrante ; la colonne ecrit la reference bien
%  formee — et c'est en verifiant ce point qu'on a decouvert que la
%  colonne tamoule, elle, avait recopie la coquille tout en annoncant
%  le contraire dans son en-tete. Corrige la, en meme temps qu'ecrit
%  ici.
%
%  LE FLEAU (40) : L'OURDOU FAIT LA SIXIEME COLONNE SANS. L'arabe
%  egyptien ecrivait عصي الدراس, le marathe मळणीच्या काठ्या, le
%  telougou నూర్పిడి కర్రలు, le tamoul போரடிக்கும் கழிகள் — quatre
%  periphrases pour le meme trou. Le Pendjab bat au pied des bœufs
%  ou au traineau, comme le Deccan et comme l'Egypte : il n'y a pas
%  d'outil a deux batons relies par une laniere, donc pas de mot. La
%  colonne ecrit گہائی کی چھڑیاں. Le coreen reste seul a avoir le mot
%  ET l'outil, 도리깨. Six colonnes sans, une avec : le trou est dans
%  la CHOSE, et la sixieme langue ne fait que confirmer ou passe la
%  ligne.
%
%  MAIS LES OISEAUX SONT TOUS LA, ET AVEC DE VIEUX NOMS. ممولا la
%  bergeronnette (19), چنڈول l'alouette (41), ابابیل l'hirondelle
%  (56), چاہا la becasse (61), چڑیا le moineau (16). Le telougou
%  avait declare trois de ces quatre non resolus ; le tamoul les
%  avait retrouves ; l'ourdou les a tous, et چاہا comme چنڈول sont
%  des mots de chasse, pas de dictionnaire. Le releve confirme ce que
%  la colonne tamoule laissait deja voir : ce sont les PLANTES qui
%  manquent aux langues d'Asie du Sud, pas les oiseaux.
%
%  ET LES FLEURS DES CHAMPS SONT PERSANES, UNE A UNE. گل گندم le
%  bleuet (1) — « la fleur du ble », qui nomme la plante par le champ
%  ou elle pousse et vaut mieux que la couleur du francais ;
%  گل انگشتانہ la digitale (3) — « la fleur en doigtier », le calque
%  persan exact du latin digitalis ; گل بہار la paquerette ; لالہ le
%  coquelicot (2) ; نیلوفر le nenuphar (18), qui est le meme mot que
%  le francais. La ou le tamoul avait du emprunter கார்ன்ஃப்ளவர் et
%  டிஜிட்டாலிஸ், l'ourdou tire tout de son registre savant : le
%  persan avait deja nomme ces plantes-la parce qu'elles poussent en
%  Iran. La coupure ne suit pas la famille de langues, elle suit la
%  carte — encore.
%
%  DEUX TROUS, ET ILS SE RESSEMBLENT. Le tilleul (35) n'a pas de nom
%  ourdou : l'arbre pousse dans l'Himalaya, pas dans la plaine, et le
%  mot n'est pas descendu avec lui. La colonne emprunte donc
%  لنڈن کا درخت — avec le ڈ retroflexe, qui le separe a l'ŒIL de
%  لندن, Londres, mais pas a l'oreille. C'est le premier emprunt de
%  tout le releve qui se distingue d'un autre mot par la seule forme
%  d'une lettre. Le chardonneret (36) est l'autre trou : la colonne
%  ecrit گولڈ فنچ, alors que le serin (35), lui, a son nom ourdou,
%  کنیری. Deux oiseaux de cage europeens, et un seul des deux a
%  atteint les cages indiennes — le nom suit l'oiseau, pas l'espece.
%
%  ET LE VIN RESTE DU VIN. La vendange, le pressoir, la cuve, le
%  tonnelier, le cidre et la biere n'ont rien d'ourdou, et la
%  tentation de les remplacer etait la meme qu'en tamoul. La colonne
%  les nomme : انگور, کولھو, بڑا ٹب, پیپا ساز, سیب کی شراب, بیئر, et
%  شراب pour le vin, sans detour. Septieme refus de ce genre, et le
%  premier dans une langue dont la plupart des lecteurs ne boivent
%  pas de vin. On modernise la langue, jamais les choses.
%
%  UNE TROUVAILLE QUI NE COUTE RIEN : پچھوا, le vent d'ouest (54).
%  L'ourdou du nord de l'Inde nomme les vents d'un seul mot —
%  پُروا celui d'est, پچھوا celui d'ouest — la ou l'ido et le
%  francais ont besoin de deux. C'est le contraire du fleau : ici
%  c'est la langue qui est plus courte que la chose.
%
%  VINGT-QUATRE PAIRES, DIX-SEPT RETOURNEMENTS. Le tableau alterne le
%  recit et l'inventaire : la moisson enumere des outils, l'orage et
%  la peche attachent. Meme alternance qu'au tableau 7, et le compte
%  monte parce que la vendange, au bloc c2-04-1, pose neuf renvois de
%  suite sur un seul tonnelier.
% ===================================================================

%%K t08-noto-1 noto
\VUnotes{143.2mm}{%
\textit{(*) کیونکہ یہ لفظ دقیق نہیں : کیا یہ سن کی کٹائی ہے، انگور کی}
\textit{کٹائی ہے، سیب کی کٹائی ہے؟ شاید Mondo XIV، ص۔ 242 میں تجویز کیا}
\textit{ہوا «} \VUgras{mesono} \textit{» استعمال کرنا بہتر ہو۔ مگر اب}
\textit{تک اُس نئے مادّے کو اکادمی نے قبول نہیں کیا ؛ اِدو کے مقررہ}
\textit{قواعد کے مطابق وہ بحث کا منتظر ہے۔}%
}
%%K t08-apar-1 apar
\VUsaut{5.0mm}
\VUtitre{15.0pt}{60}{جدول نمبر 8}
\VUsaut{3.0mm}
\VUcentre{13.2pt}{90}{\textit{پہلا منظر۔}}
\VUsaut{3.0mm}
\VUpk{10.2pt}{40}{فصل کی کٹائی (*)۔}
\VUsaut{4.0mm}
\VUpk{10.2pt}{40}{دیہات کے مناظر۔}
\VUsaut{3.0mm}

%%K t08-c1-01-1 p
1. --- \VUgras{گرمیوں} کے ساتھ دیہات کا منظر ایک بار پھر بدل گیا۔
نالی کے سپرد کیا ہوا بیج پھوٹا ؛ سبز ننھا پودا زمین سے نکلا اور بالی
لے آیا۔ دانہ بنا، پھر پکا، اور اب اُسے کاٹنے کے سوا کچھ باقی نہیں۔

%%K t08-c1-02-1 p
2. --- \VUgras{دیہاتی پھول} کھل چکے ہیں۔ \VUgras{گل گندم}
\textsuperscript{(1)}، \VUgras{لالہ کے ننھے پھول}
\textsuperscript{(2)} اور \VUgras{گل انگشتانہ} \textsuperscript{(3)}
گندم کے کھیتوں کے یکساں رنگ میں اپنی تازہ \VUgras{پنکھڑیوں} کا
شاداب تضاد ملا دیتے ہیں۔

%%K t08-c1-03-1 p
3. --- پھل دار درخت پتوں سے بھر گئے ؛ پھل پک گیا۔ چھوٹے
\VUgras{گیلاس کے درخت} \textsuperscript{(4)} پر خوبصورت \VUgras{گیلاس}
\textsuperscript{(5)} لدے ہیں ؛ بچے اُنھیں بڑے شوق سے توڑتے اور کھاتے
ہیں۔

%%K t08-c1-04-1 p
4. --- اب سے مویشی سارا دن کھلی ہوا میں گزار سکتے ہیں۔

%%K t08-c1-04-2 p
\VUgras{میمنے} \textsuperscript{(6)} بڑے ہو گئے اور گھنی اُون والی
خوبصورت \VUgras{بھیڑیں} \textsuperscript{(7)} اور خوبصورت
\VUgras{مینڈھے} \textsuperscript{(8)} بن گئے۔ وہ اطمینان سے
\VUgras{چراگاہ} \textsuperscript{(9)} کی \VUgras{گھاس} چرتے ہیں ؛ وہ
چراگاہ \VUgras{گل بہار} سے رنگی ہوئی ہے۔

%%K t08-c1-04-3 p
جلد ہی \VUgras{اُن} جانوروں کے بال کترے جائیں گے تاکہ اُن کی
\VUgras{اُون} ملے ؛ اُسی سے ہمارے کپڑوں کا کھدر بنتا ہے۔

%%K t08-tit-2 sub
\VUsaut{5.0mm}
\VUpk{10.2pt}{40}{گڈریا۔}
\VUsaut{3.4mm}

%%K t08-c1-05-1 p
5. --- \VUgras{گڈریا} \textsuperscript{(10)} اُن کی زیادہ پروا کرتا نہیں
لگتا۔ اُسے صرف اُس طوفان کی فکر ہے جو افق پر سے گزر رہا ہے ؛ اور
\VUgras{بھیڑوں کا نگہبان} \textsuperscript{(13)}، جو اپنا
\VUgras{تونبا} \textsuperscript{(14)} ہونٹوں سے لگاتا ہے، اُس سے بھی
زیادہ بے فکر لگتا ہے۔

%%K t08-c1-06-1 p
6. --- گرمی دم گھونٹ رہی ہے ؛ کیونکہ \VUgras{کتا} \textsuperscript{(15)}
بھی اپنی پیاس بجھانے آتا ہے ؛ وہ \VUgras{ندی} \textsuperscript{(16)} کا
ٹھنڈا پانی چاٹتا ہے اور کنارے کی گھاس میں دبکے ہوئے بے چارے
\VUgras{ننھے مینڈک} \textsuperscript{(17)} کو ڈرا دیتا ہے۔ وہ صاف پانی
میں کود جاتا ہے ؛ اُس پانی میں \VUgras{نیلوفر} \textsuperscript{(18)}
کھلتے ہیں۔

%%K t08-c1-07-1 p
7. --- \VUgras{ممولا} \textsuperscript{(19)} اُس ننھے \VUgras{چشمے}
\textsuperscript{(20)} کے پاس نہاتا ہے جو دو چٹانوں کے بیچ سے پھوٹتا
ہے، اور \VUgras{کانٹے دار جھاڑیوں} \textsuperscript{(21)} اور
\VUgras{زعرور کی جھاڑیوں} \textsuperscript{(22)} کے نیچے چھپ جاتا ہے۔

%%K t08-tit-3 sub
\VUsaut{5.0mm}
\VUpk{10.2pt}{40}{فصل کی کٹائی۔}
\VUsaut{3.4mm}

%%K t08-c1-08-1 p
8. --- دیہات کے بچوں کے لیے گندم کی کٹائی بہت دل چسپ موسم ہے۔ وہ
خوشی خوشی \VUgras{قلابازیاں} \textsuperscript{(24)} کھاتے ہیں —
\VUgras{بھوسے کے ڈھیروں} \textsuperscript{(23)} پر۔ وہ \VUgras{فصل
کاٹنے والوں} \textsuperscript{(26)} کی مدد کرتے ہیں ؛ اور جب وہ
\VUgras{گندم کا کھیت} \textsuperscript{(27)} اپنی \VUgras{بڑی
درانتیوں} \textsuperscript{(25)} سے کاٹتے ہیں — جو \VUgras{سان}
\textsuperscript{(30)} پر خوب تیز کی ہوئی ہیں —، بچے گندم سمیٹ کر اُسے
\VUgras{پولوں} \textsuperscript{(31)} یا \VUgras{چھوٹے گٹھوں}
\textsuperscript{(32)} میں باندھتے ہیں۔

%%K t08-c1-09-1 p
9. --- کبھی کبھی اُن میں سب سے بڑوں کو \VUgras{درانتی}
\textsuperscript{(12)} سے وہ \VUgras{سنہری ڈنٹھل} \textsuperscript{(28)}
کاٹنے کی اجازت ملتی ہے جن پر دانوں سے بھری بھاری \VUgras{بالیاں}
\textsuperscript{(29)} لگی ہیں ؛ اور چھوٹے، \VUgras{مویشیوں}
\textsuperscript{(33)} کی نگرانی کرتے ہوئے — وہ مویشی \VUgras{مرغزار}
\textsuperscript{(34)} میں بڑے \VUgras{لنڈن کے درخت}
\textsuperscript{(35)} کے سائے میں چرتے ہیں —، اپنے \VUgras{جال}
\textsuperscript{(36)} سے خوبصورت \VUgras{تتلیاں} پکڑتے ہیں۔

%%K t08-c1-09-2 p
کچھ تو \VUgras{چھکڑے} \textsuperscript{(37)} پر چڑھ کر وہ پولے ترتیب سے
رکھتے ہیں جو اُنھیں \VUgras{کانٹے} \textsuperscript{(38)} سے پکڑائے
جاتے ہیں۔

%%K t08-c1-10-1 p
10. --- سارے \VUgras{میدان} \textsuperscript{(39)} میں چہل پہل ہے ؛ اُسی
میدان سے \VUgras{چنڈول} \textsuperscript{(41)} اُڑتے ہیں۔ سب فصل کی
کٹائی میں لگے ہیں۔ مگر جہاں غریب لوگ پرانے اوزار ہی برتتے رہتے ہیں —
\VUgras{بڑی درانتیاں}، \VUgras{درانتیاں} اور \VUgras{گہائی کی
چھڑیاں} \textsuperscript{(40)} —، وہاں امیر زمین دار اپنی گندم یا اپنا
\VUgras{چاودار} \textsuperscript{(43)} \VUgras{فصل کاٹنے کی مشین}
\textsuperscript{(42)} سے کٹواتے ہیں۔ پھر، پولوں کو کھلیان تک لے جانے
کی ضرورت کے بغیر، وہیں \VUgras{پولوں کے بڑے ڈھیر}
\textsuperscript{(44)} لگا دیے جاتے ہیں۔

%%K t08-c1-11-1 p
11. --- تب \VUgras{گہائی کی مشین} \textsuperscript{(47)} لائی جاتی ہے ؛
اُسے \VUgras{بھاپ کا انجن} \textsuperscript{(45)} ایک \VUgras{پٹّے}
\textsuperscript{(46)} کے ذریعے چلاتا ہے ؛ اور یوں دانہ
\VUgras{بھوسے} سے الگ ہو جاتا ہے، بغیر اُس کھیت کو چھوڑے جس میں وہ
بویا گیا تھا۔

%%K t08-tit-4 sub
\VUsaut{5.0mm}
\VUpk{10.2pt}{40}{طوفان۔}
\VUsaut{3.4mm}

%%K t08-c1-12-1 p
12. --- کٹائی کے موسم میں اکثر \VUgras{زبردست طوفان}
\textsuperscript{(53)} آتے ہیں۔ پہلے دور سے \VUgras{گرج} کی گڑگڑاہٹ
سنائی دیتی ہے ؛ \VUgras{پچھوا آندھی} \textsuperscript{(54)} چلنے لگتی ہے
اور ہانپتی ہوئی \VUgras{چھوٹے جنگل} \textsuperscript{(49)} کے موٹے سے
موٹے درختوں کو جھکا دیتی ہے۔ کالے \VUgras{بادل} \textsuperscript{(56)}
آسمان پر چڑھ آتے ہیں ؛ اُنھیں \VUgras{بجلی} \textsuperscript{(55)} کی
چندھیا دینے والی ٹیڑھی لکیریں چیرتی ہیں۔ \VUgras{بارش} اور کبھی کبھی
\VUgras{اولے} گرنے لگتے ہیں، اور کٹنے کو تیار فصل کو کچل دیتے ہیں۔

%%K t08-c1-13-1 p
13. --- اکثر بجلی کسی عمارت کو جلا دیتی ہے۔ ایسے ہی پچھلے سال
\VUgras{پون چکی} \textsuperscript{(50)} کی چھت راکھ ہو گئی اور اُس کے دو
\VUgras{پنکھے} \textsuperscript{(51)} چور چور ہو گئے۔

%%K t08-c1-14-1 p
14. --- کچھ گھنٹوں بعد سکون لوٹ آتا ہے۔ افق پر چمکتی \VUgras{دھنک}
\textsuperscript{(52)} نمودار ہوتی ہے اور فطرت اپنی وہ زندگی پھر شروع
کر دیتی ہے جو چند لمحوں کے لیے رُک گئی تھی۔ دیہاتی، جنھوں نے
\VUgras{جھونپڑیوں} \textsuperscript{(48)} میں پناہ لی تھی، پھر کام پر لگ
جاتے ہیں اور ہمت سے اُس نقصان کی تلافی کی کوشش کرتے ہیں جو اُنھیں
ابھی اٹھانا پڑا۔

%%K t08-tit-5 sub
\VUsaut{6.0mm}
\VUcentre{13.2pt}{90}{\textit{دوسرا منظر۔}}
\VUsaut{3.6mm}

%%K t08-tit-6 sub
\VUpk{10.2pt}{40}{شکار۔ --- انگور کی کٹائی۔}
\VUsaut{1.4mm}
\VUpk{10.2pt}{40}{مچھلی کا شکار۔}
\VUsaut{4.0mm}

%%K t08-c2-01-1 p
1. --- علاقے کے مطابق \VUgras{اگست} کے آخر یا \VUgras{ستمبر} کے وسط
میں شکار کا موسم کھلتا ہے۔ سورج کی پہلی کرنوں نے ابھی \VUgras{ننھی
چھپکلیوں} \textsuperscript{(2)} کو اُن کے بل سے نکالا ہی تھا کہ
\VUgras{شکاری} \textsuperscript{(5)} اپنے \VUgras{کتوں}
\textsuperscript{(4)} کے ساتھ نکل پڑا۔

%%K t08-c2-02-1 p
2. --- اُس نے موٹے کھدر کا اپنا مضبوط \VUgras{شکاری لباس}
\textsuperscript{(9)} پہنا، پاؤں میں چمڑے کے \VUgras{پنڈلی پوش} چڑھائے
— اُن کے سہارے وہ اُس گیلی گھاس میں چل سکے گا جہاں \VUgras{کھردرے
مینڈک} \textsuperscript{(1)} چھپے رہتے ہیں —، اپنی \VUgras{بندوق}
\textsuperscript{(6)} اور اپنا \VUgras{شکار کا تھیلا}
\textsuperscript{(10)} اٹھایا، اور یہ چلا۔

%%K t08-c2-03-1 p
3. --- تعاقب کا جوش اُسے انگور کے ایک باغ کے بیچ لے آتا ہے جہاں
انگور کی کٹائی زوروں پر ہے۔ انگور کے کاشتکار کے بچے، جو عام طور پر
سڑک پر بکریاں چراتے ہیں، آج اُنھیں جہاں چاہیں چرنے کو چھوڑ دیتے ہیں،
اور ایک بوڑھے \VUgras{بکرے} \textsuperscript{(3)} کو کھونٹے سے باندھنے
کے بعد — جو اپنی آزادی سے فائدہ اٹھا کر بھاگ نکلنے میں ہرگز نہ چوکتا
—، وہ بھی اپنے حصے کی خوشی لے لیتے ہیں۔ ایک \VUgras{انگور توڑنے کی
ٹوکری} \textsuperscript{(12)} سے انگور کا ایک گچھا اٹھاتا ہے اور
\VUgras{رس} \textsuperscript{(14)} \VUgras{کٹورے} \textsuperscript{(13)}
میں نچوڑ کر مست بنانے میں دل بہلاتا ہے (وہ شراب جس میں ابھی خمیر
نہیں اٹھا)۔ دوسرا \VUgras{پتوں کا تاج} \textsuperscript{(15)} سر پر
رکھتا ہے، جو اُس نے ابھی گوندھا ہے۔

اُن کے پاس ہی بے شرم \VUgras{چڑیاں} \textsuperscript{(16)} انگور کے
دانے چگنے کولھو کے سامنے تک آ جاتی ہیں۔

%%K t08-c2-04-1 p
4. --- جب بچے کھیل رہے ہیں، مرد کام کر رہے ہیں۔ \VUgras{انگور توڑنے
والے} \textsuperscript{(18)} \VUgras{پیٹھ کی ٹوکریوں}
\textsuperscript{(17)} میں انگور تہ خانے تک لاتے ہیں ؛ \VUgras{پیپا
ساز} \textsuperscript{(19)} \VUgras{پیپوں} \textsuperscript{(20)} کی
مرمت کرتا ہے، دیکھتا ہے کہ \VUgras{پیپے کے تختے}
\textsuperscript{(22)} اور \VUgras{پیندے} \textsuperscript{(21)} ٹھیک
حالت میں ہیں یا نہیں، اور ٹوٹے \VUgras{لوہے کے حلقے}
\textsuperscript{(23)} بدل دیتا ہے۔ وہی \VUgras{بڑے ٹب}
\textsuperscript{(25)} بھی بناتا ہے ؛ اُن میں \VUgras{انگور}
\textsuperscript{(26)} ڈال کر خمیر اٹھایا جاتا ہے۔

%%K t08-c2-05-1 p
5. --- خمیر اٹھ جانے کے بعد شراب نتھار لی جاتی ہے اور بچا ہوا گودا
\VUgras{کولھو} \textsuperscript{(28)} پر رکھا جاتا ہے ؛ \VUgras{پیچ}
\textsuperscript{(29)} کو رفتہ رفتہ کس کر رس نچوڑا جاتا ہے ؛ وہ رس
\VUgras{چھوٹے ٹب} \textsuperscript{(32)} میں بہتا ہے ؛ اور یوں پھر
\VUgras{شراب} \textsuperscript{(31)} حاصل ہوتی ہے۔

%%K t08-tit-8 sub
\VUsaut{6.0mm}
\VUpk{10.2pt}{40}{بیئر اور سیب کی شراب۔}
\VUsaut{3.4mm}

%%K t08-c2-06-1 p
6. --- \VUgras{تہ خانے} \textsuperscript{(27)} کی دیوار پر
\VUgras{ہاپس} \textsuperscript{(30)} کی ایک شاخ چڑھتی ہے۔ وہاں یہ صرف
سجاوٹ کا پودا ہے، مگر کچھ ملکوں میں، جہاں انگور کی بیل بہت کم ہے،
ہاپس \VUgras{بیئر} بنانے کے لیے اُگایا جاتا ہے۔

%%K t08-c2-07-1 p
7. --- چھوٹا \VUgras{سیب کا درخت} \textsuperscript{(33)} سیبوں سے لدا
ہے ؛ وہ سیب پک جانے پر بہترین \VUgras{سیب کی شراب} دیں گے۔ اپنے
\VUgras{پنجرے} \textsuperscript{(34)} میں \VUgras{کنیری}
\textsuperscript{(35)} اور \VUgras{گولڈ فنچ} \textsuperscript{(36)} اُن
چڑیوں کو حسد بھری آنکھوں سے دیکھتے ہیں جو آزاد پھدکتی پھرتی ہیں۔

%%K t08-c2-08-1 p
8. --- حد نگاہ تک، ٹیلوں کی ڈھلان پر \VUgras{انگور کے کھونٹے}
\textsuperscript{(40)} قطار در قطار لگے ہیں ؛ وہ نہایت خوبصورت
\VUgras{انگور کی بیلیں} \textsuperscript{(39)} تھامے ہوئے ہیں ؛ اُن
بیلوں پر بھاری \VUgras{گچھے} لدے ہیں۔

%%K t08-c2-08-2 p
سارا سال \VUgras{انگور کا کاشتکار} \textsuperscript{(42)} اپنے
\VUgras{انگور کے باغ} \textsuperscript{(38)} کی دیکھ بھال میں محنت کرتا
رہا، اور اب اُسے اُس محنت کا صلہ خوب فصل کی صورت میں ملتا ہے۔
\VUgras{انگور توڑنے والیاں} \textsuperscript{(41)} اُس چھکڑے پر رکھے
ٹب بھرتی ہیں جسے \VUgras{خچریاں} \textsuperscript{(43)} کھینچتی ہیں،
اور سب خوش ہیں۔

%%K t08-tit-9 sub
\VUsaut{5.0mm}
\VUpk{10.2pt}{40}{مچھلی کا شکار۔}
\VUsaut{3.4mm}

%%K t08-c2-09-1 p
9. --- یہی حال \VUgras{بنسی سے مچھلی پکڑنے والوں} \textsuperscript{(44)}
کا ہے ؛ وہ صبح ہی سے اپنی \VUgras{بنسی کی چھڑیاں}
\textsuperscript{(45)} لے کر پانی کے کنارے آ بیٹھے ہیں۔

%%K t08-c2-09-2 p
\VUgras{مچھلی} \textsuperscript{(47)} \VUgras{کانٹے} پر جھپٹتی ہے، جوں
ہی وہ اپنی \VUgras{ڈور} \textsuperscript{(46)} پانی میں ڈالتے ہیں ؛ اور
اُنھیں \VUgras{چارا} بدلنے کی مہلت بھی مشکل سے ملتی ہے کہ ڈور کے سرے
پر نیا شکار پھڑکنے لگتا ہے۔

%%K t08-c2-09-3 p
پھر بھی \VUgras{جال سے مچھلی پکڑنے والا} \textsuperscript{(48)} زیادہ
مچھلیاں پکڑتا ہے ؛ وہ اپنی \VUgras{کشتی} \textsuperscript{(49)} سے
\VUgras{پھینکنے کا جال} \VUgras{دریا} \textsuperscript{(50)} کے بیچوں
بیچ پھینکتا ہے۔

%%K t08-c2-10-1 p
10. --- خزاں اچھی سیروں کا موسم ہے : پیدل، \VUgras{گھوڑے}
\textsuperscript{(52)} پر، \VUgras{سائیکل} \textsuperscript{(53)} پر،
\VUgras{موٹر گاڑی} \textsuperscript{(54)} میں۔ \VUgras{سڑکیں}
\textsuperscript{(51)} صاف ہیں اور لوگ آخری خوشگوار دنوں سے فائدہ
اٹھاتے ہیں، کیونکہ یہ رہیں \VUgras{ابابیلیں} \textsuperscript{(56)} جو
چھتوں پر جمع ہو رہی ہیں تاکہ پورے پروں سے زیادہ گرم ملکوں کی طرف اُڑ
جائیں۔ بوڑھے \VUgras{اخروٹ کے درخت} \textsuperscript{(57)} کے پتے زرد
ہو رہے ہیں، \VUgras{اخروٹ} گر رہے ہیں، اور \VUgras{گلدستے}
\textsuperscript{(58)} کی طرح جھنڈ باندھے \VUgras{ٹیلے}
\textsuperscript{(59)} پر کھڑے اونچے \VUgras{درخت} جلد ہی بے پتے ہو
جائیں گے۔

%%K t08-c2-11-1 p
11. --- \VUgras{سورج} \textsuperscript{(60)} دیر سے نکلتا ہے، دن چھوٹے
ہو رہے ہیں، \VUgras{صبح} کو خاصی چبھتی ٹھنڈ محسوس ہونے لگی ہے، اور یہ
رہیں \VUgras{چاہوں} \textsuperscript{(61)} کی ڈاریں جو ہماری بے مہمان
نواز آب و ہوا چھوڑ رہی ہیں۔

\VUsaut{6.0mm}
\VUfilet{22mm}
